\section{Dataset Overview}

The dataset covers daily observations from April 2018 to April 2023, a period
that spans the late-cycle expansion of 2018--2019, the COVID-19 pandemic shock
of early 2020, the subsequent monetary and fiscal policy response, the
post-pandemic inflation surge, and the aggressive central-bank tightening
campaigns of 2022--2023. Below we describe each series and its behaviour over
the sample.

\subsection*{Equity indices}

\begin{itemize}
  \item \textbf{S\&P 500 (\texttt{S\&P}).}
    The US large-cap equity benchmark gained approximately 53\% over the
    sample, but the path was far from linear. A sharp sell-off in March 2020,
    triggered by global lockdowns, saw the index fall to 2\,237 before
    substantial Fed easing and fiscal stimulus drove a rapid recovery. It
    reached an all-time high near 4\,800 in January 2022, then retreated
    roughly 20\% during the rate-hiking cycle before partially recovering into
    early 2023.

  \item \textbf{EuroStoxx 600 (\texttt{ESTX600}).}
    The broad European equity benchmark rose approximately 22\% over the
    period, following a broadly similar pattern---a crash to 280 in March 2020,
    a strong recovery, and a high near 494 in early 2022. However, its
    performance lagged US equities, reflecting Europe's greater exposure to the
    energy crisis triggered by the Russia--Ukraine conflict and relatively
    slower earnings growth.
\end{itemize}

\subsection*{Foreign exchange}

\begin{itemize}
  \item \textbf{EURUSD spot rate (\texttt{EURUSD}).}
    The euro depreciated approximately 12\% against the dollar over the sample,
    driven by the Fed tightening faster and more aggressively than the ECB,
    compounded by the European energy crisis. The pair briefly broke below
    parity in September 2022, reaching 0.9594---a level not observed since
    2002---before recovering as the ECB accelerated its own hiking cycle.

  \item \textbf{EURUSD 3-month FX basis swap (\texttt{EURUSD3mBS}).}
    This series measures the cost of swapping EUR for USD funding over three
    months, quoted in basis points. On a point-to-point basis it ended close
    to where it started, but exhibited extreme volatility in between. The most
    notable move was a collapse to --81\,bps in March 2020, reflecting acute
    dollar funding stress as global institutions sought USD liquidity.
    The Fed's emergency swap lines quickly normalised the basis. It
    subsequently swung into positive territory, peaking near +65\,bps, as the
    ECB's rate hikes made EUR funding relatively more expensive.
\end{itemize}

\subsection*{Government bond yields and yield curves}

\begin{itemize}
  \item \textbf{10-year US Treasury yield (\texttt{UST10}).}
    The 10-year yield experienced a significant round-trip: it fell to a
    historic low of 0.51\% in August 2020 as the Fed cut rates to zero and
    launched large-scale asset purchases, then rose to 4.24\% in
    October 2022, driven by persistent inflation and aggressive tightening.
    It ended the sample moderately higher than its starting level.

  \item \textbf{10-year Bund yield (\texttt{DEM10}).}
    The German 10-year benchmark fell into deeply negative territory
    (--0.86\%) during 2019--2020, reflecting the ECB's negative-rate policy
    and safe-haven demand during COVID. The yield turned positive in early 2022
    and peaked at 2.75\% as the ECB pivoted to rate hikes, ending the sample
    materially above its starting level.

  \item \textbf{2-year US Treasury yield (\texttt{UST2}).}
    The 2-year yield bottomed at 0.10\% in 2020--2021 when the Fed held rates
    at the zero lower bound, then rose to 5.07\% as the market priced in
    cumulative Fed hikes of over 500\,bps. The short end led the move higher
    because 2-year rates are more tightly anchored to policy-rate expectations,
    ending the sample significantly above its initial level.

  \item \textbf{2-year Schatz yield (\texttt{DEM2}).}
    The German 2-year yield traded as low as --1.01\%, deep in negative
    territory while the ECB deposit rate sat at --0.50\%. When the ECB began
    hiking in July 2022 it climbed rapidly, peaking at 3.33\%---reflecting
    one of the fastest policy pivots in ECB history---and ending the sample
    well above its starting level.

  \item \textbf{2s--10s US Treasury curve (\texttt{USYC}).}
    The 10-year minus 2-year spread steepened to +158\,bps in early 2021 as
    the market anticipated post-COVID reflation while the Fed held the front
    end at zero. It then inverted sharply, reaching --108\,bps, as the
    2-year yield outpaced the 10-year during the hiking cycle---a classic
    recession warning signal.

  \item \textbf{2s--10s German curve (\texttt{DEYC}).}
    The equivalent German spread was persistently steep during the ECB's
    negative-rate era, then flattened and inverted as ECB tightening pulled the
    short end higher faster than the long end. The transition from a positively
    sloped to an inverted curve mirrored the US dynamic with a lag.
\end{itemize}

\subsection*{Overnight rates}

\begin{itemize}
  \item \textbf{SOFR (\texttt{SOFR}).}
    The Secured Overnight Financing Rate, the Fed's preferred short-term
    benchmark, fell to near zero (0.01\%) in March 2020 following emergency
    inter-meeting rate cuts, then rose steadily to 4.80\% by the end of the
    sample, closely tracking the effective federal funds rate through the full
    hiking cycle.

  \item \textbf{Euro short-term rate (\texttt{ESTR}).}
    The ECB's overnight benchmark was pinned around --0.55\% through 2021,
    reflecting the negative deposit facility rate. It began rising in mid-2022
    and ended at 2.90\%, mirroring the ECB's cumulative 350\,bps of
    tightening.
\end{itemize}

\subsection*{Credit default swap indices}

\begin{itemize}
  \item \textbf{CDX IG 5-year (\texttt{CDX IG}).}
    The North American investment-grade CDS index ended the sample modestly
    wider than its starting level. It rose to 152\,bps during the
    March 2020 liquidity crisis as corporate default concerns intensified. The
    Federal Reserve's decision to backstop the corporate bond market via the
    Secondary Market Corporate Credit Facility (SMCCF) rapidly compressed
    spreads back toward pre-crisis levels. A secondary widening occurred in
    mid-2022 amid recession fears, though materially smaller than the COVID
    episode.

  \item \textbf{CDX HY 5-year (\texttt{CDX HY}).}
    The high-yield counterpart widened to 871\,bps in March 2020, nearly
    tripling from pre-crisis levels, as high-yield issuers were seen as most
    vulnerable to a prolonged shutdown. The recovery was swift but HY spreads
    remained wider than their pre-COVID tights throughout the sample, ending
    materially above the starting level.

  \item \textbf{iTraxx Main IG 5-year (\texttt{ITRX IG}).}
    The European investment-grade CDS index mirrored CDX IG's COVID spike
    (peaking at 139\,bps) and subsequent compression, though European spreads
    were additionally affected by the 2022 energy crisis and sovereign-spread
    concerns in peripheral Europe. It ended the sample moderately wider.

  \item \textbf{iTraxx Crossover 5-year (\texttt{ITRX XO}).}
    The European sub-investment-grade/crossover index peaked at 708\,bps in
    March 2020. Crossover spreads remained structurally wider than pre-COVID
    levels throughout the sample, partly reflecting lingering uncertainty
    around European energy security and the corporate impact of higher rates.
\end{itemize}

\subsection*{Commodities}

\begin{itemize}
  \item \textbf{Brent crude oil (\texttt{CO}).}
    Brent crude gained approximately 49\% over the sample but with significant
    intra-period volatility. It fell sharply in early 2020 as lockdowns
    destroyed transport fuel demand, then rose above \$100 in mid-2022
    following Russia's invasion of Ukraine and associated supply disruptions.
    Prices moderated into 2023 as demand concerns from the global slowdown
    partially offset tight supply.
\end{itemize}

\subsection*{Inflation breakevens}

\begin{itemize}
  \item \textbf{US 5-year inflation breakeven (\texttt{US5BEI}).}
    On a point-to-point basis the breakeven ended near its starting level,
    but the intra-sample range was notable: it fell to 0.18\% in March 2020
    as markets priced in a deflationary demand shock, then rose to 3.73\% by
    mid-2022 as realised CPI inflation exceeded 9\%. The magnitude of this
    round-trip was among the largest observed historically, reflecting a rapid
    transition from deflation expectations to sustained above-target inflation.

  \item \textbf{French 5-year inflation breakeven (\texttt{FR5BEI}).}
    This series is excluded from all quantitative analyses due to insufficient
    data coverage (85\% of observations missing). A detailed justification is
    provided in the sensitivity document (Section~2).
\end{itemize}

\subsection*{Volatility}

\begin{itemize}
  \item \textbf{VIX (\texttt{VIX}).}
    The CBOE implied-volatility index ended the sample near its starting
    level, consistent with the VIX's well-known mean-reverting properties. The
    defining event was a rise to 82.7 in March 2020---its
    highest reading since the 2008 financial crisis---as markets experienced
    extreme dislocations. It subsequently mean-reverted, with periodic
    flare-ups during the 2022 sell-off, before settling back near long-run
    averages.
\end{itemize}
